% SOURCE_STATEMENT: "Problem 28. Consider negative energy three-body orbits which are unbound in both direc-tions of time. The two Jacobi vectors of such an orbit then asymptote to Keplerian orbits in both the distant past and the distant future, and so associated to such an orbit we have pair of Kepler elements in the past (one elliptic the other hyperbolic) and another pair in the future. The \"direct scattering\" problem is: which pair of Kepler elements can be connected to each other in this way? (Richard Montgomery)"
% Preparation and final self-audit: CONTRIBUTING.md prompts C and D; source/audit artifacts retained in the local physics_statements_2026_10_11 batch.
% !TeX program = lualatex
\documentclass[11pt]{article}
\usepackage[margin=25mm]{geometry}
\usepackage{fontspec}
\setmainfont{Latin Modern Roman}
\usepackage{amsmath,amssymb,mathtools}
\usepackage{unicode-math}
\setmathfont{Latin Modern Math}
\usepackage{xurl,hyperref}
\hypersetup{colorlinks=true,urlcolor=blue}
\setcounter{secnumdepth}{0}
\setlength{\parindent}{0pt}
\setlength{\parskip}{5pt}
\setlength{\emergencystretch}{3em}
\newcommand{\sref}[2]{\hyperlink{source:#1:#2}{[S#2]}}
\newcommand{\cataloguescope}[1]{\paragraph{Recorded target.}#1}
\newcommand{\R}{\mathbb R}
\newcommand{\C}{\mathbb C}
\newcommand{\N}{\mathbb N}
\newcommand{\Z}{\mathbb Z}
\newcommand{\Q}{\mathbb Q}
\newcommand{\F}{\mathbb F}
\newcommand{\eps}{\varepsilon}
\begin{document}
\section{A feasibility sieve for negative-energy binary-single direct scattering}\label{statement}

\subsection{Definitions and mathematical statement}

\paragraph{Shared notation.}$\mathbb N=\{1,2,\ldots\}$; the standard number fields have their usual meanings. All additional conventions are stated locally.


\paragraph{Definitions and notation.}In the Newtonian $n$-body problem, positions $q_i\in\mathbb R^d$ and masses $m_i>0$ obey $m_i\ddot q_i=\nabla_{q_i}U$, where $U(q)=\sum_{i<j}m_im_j/|q_i-q_j|$ and the gravitational constant is normalized to one. Here $T=\tfrac12\sum_i m_i|\dot q_i|^2$, $E=T-U$, and a collision has $q_i=q_j$. The center of mass can be fixed at zero; moment of inertia is $I=\sum_i m_i|q_i|^2$. For a binary plus third body, Jacobi coordinates describe internal binary separation and the third body’s position relative to the binary center of mass. Kepler elements encode an asymptotic two-body conic, including energy, angular momentum, orientation and phase/time conventions. Negative total energy can coexist with an escaping third body when the binary remains bound.

\paragraph{Catalogue wording retained for traceability.}
Consider three-body orbits with negative total energy that are unbound in both time directions. Their Jacobi vectors approach a bound elliptic binary and a hyperbolic escaping orbit in each distant past/future description. Which pairs of incoming and outgoing Kepler elements can be connected by one such orbit?
\paragraph{Scope and source qualification.}Energy and angular-momentum conservation give necessary constraints rather than a full scattering characterization. The source’s asymptotic Kepler description is preserved; phase/time conventions and handling of channel exchange must be specified. The catalogue title/status was paired with a machine-generated research attempt; that attempt is not peer-reviewed evidence of settlement. The statement covers the original source questions rather than only the title’s subcase.
\subsection{Short English statement}
Characterize the incoming and outgoing Kepler data connected by negative-energy three-body scattering.
\subsection{Sources}
\begin{itemize}
\item[\hypertarget{source:statement:1}{S1}] ULAM.AI and the UnsolvedMath Contributors, UnsolvedMath, catalogue entry. Catalogue curation under CC BY 4.0; definitions and English exposition added locally.
\url{https://huggingface.co/datasets/ulamai/UnsolvedMath}.
\item[\hypertarget{source:statement:2}{S2}] American Institute of Mathematics, Variational Methods in Celestial Mechanics; notes by Kuo-Chang Chen, version 22 June 2003, source item 28.
\url{https://aimath.org/WWN/varcelest/varcelest.pdf}.
\end{itemize}
\label{end:statement}
\end{document}
